\section{Model and information-certified variance bound}
\label{sec:model}\label{sec:certificate}

The systematic exposures that generate portfolio risk are latent, whereas the
paper observes firms' information.
The model therefore separates three roles: an observed information law, a
latent systematic-exposure law, and a maintained transmission restriction that
connects the two without equating them.
It then places all latent exposures under one coherent joint law so that
pairwise restrictions can support an $n$-asset portfolio statement.

Let $i\in\{1,\ldots,n\}$ index assets.
The observed object records the distribution of firm $i$'s information.
Formally, $X_i$ is an observable characteristic draw taking values in a
separable metric space $\mathcal X$, and its law $C_i$ is the probability
distribution of the firm's row-normalized article embeddings.
The latent object records the corresponding systematic exposure.
Let $B_i$ denote that exposure in a real Hilbert space $\mathcal H$.
The exposure model is
\begin{equation}\label{eq:p3-exposure-model}
  B_i=u_i(X_i),
  \qquad P_i=\mathcal L(B_i),
\end{equation}
where the measurable map $u_i$ captures the firm's information-to-exposure
mapping and $P_i$ is the resulting latent exposure law.
Thus $C_i$ and $P_i$ are distinct laws and need not use the same metric units.

The maintained transmission restriction supplies the link between these two
spaces.
Its common component is a measurable carrier
$t:\mathcal X\to\mathcal H$.
It is $L$-antilipschitz:
\begin{equation}\label{eq:p3-antilipschitz}
  d_{\mathcal H}(t(x),t(y))
  \ge L^{-1}d_{\mathcal X}(x,y),
  \qquad L>0.
\end{equation}
Each firm-specific map remains within a synchronous slack radius:
\begin{equation}\label{eq:p3-synchronous-slack}
  d_{\mathcal H}(u_i(x),t(x))\le\tau_i
  \quad\text{for every }x,
  \qquad \tau_i\ge0.
\end{equation}
The antilipschitz restriction prevents economically distinct information states
from collapsing into identical systematic exposures, while $\tau_i$ permits
firm-specific information-to-exposure mismatch.
Together, these restrictions translate observed information distance into a
lower bound on exposure separation rather than an equality or a covariance
estimate.

Pairwise exposure restrictions do not by themselves define portfolio risk.
To aggregate them, let $J$ be one coherent joint law for
$(B_1,\ldots,B_n)$ with coordinate marginals $P_i$.
For normalized risk weights $q_i\ge0$ satisfying $\sum_i q_i=1$, define
\[
  v_i=\mathbb E_J\lVert B_i\rVert^2,
  \qquad
  V_{\mathrm{sys}}(q)
  =\mathbb E_J\left\lVert\sum_i q_i B_i\right\rVert^2.
\]
The same $J$ governs every cross term, so the induced covariance matrix is
positive semidefinite.
This requirement rules out assembling a portfolio from mutually incompatible
pairwise optimal couplings.

With the observed and latent objects linked and joint coherence imposed, the
next subsections derive observable pairwise floors, aggregate them into a
conservative portfolio certificate, sharpen that certificate with a multi-firm
object, and finally connect systematic exposure risk to returns.
